\chapter{厚度、档位与探针的对照结果}
\label{ch:result}

数值比较要回答三件事：开环菜单在什么厚度上失败，腔温闭环闭合的是空气还是冷点，以及 40 s 探针有没有把内部热阻的方差降下来。\(\alpha_m=0.85\) 是策略比较的零点，对应恒温冷点上 \(F=3.794\,\mathrm{min}\)。表中每一行都是同一数值核在因子格子上的一条样本；读表时先按厚度拆开，再比较策略名称，因为食品名、初始含水和档位名称的效应都小于厚度。交叉表同时也是产品说明书里那张“智能菜单效果”对照的展开写法。说明书通常只给一张成功照片：表面已经上色，切开后中心不再是粉色。格子把同一张照片拆成 19440 行，覆盖三种厚度、两种初温、三种含水、三种非均质和四种策略，每个格子 30 个种子。成功照片对应的是 30 mm、初温 \(8^\circ\mathrm{C}\)、对偶、组合那一条，停机 1866 s，\(e_{\alpha}=+0.145\)；全表中位却是 \(-0.313\)，因为格子里混有 50 mm 和 \(5^\circ\mathrm{C}\)。把演示页的默认条件当成全表中位，会把一次成功读成“对偶已经越过设计线”。

\section{对照试验的格子设计}

主对照 S1 的因子是三种食品、三种厚度、两种初温、三种含水、三种非均质和四个策略，每个格子 30 个种子，合计
\[
3\times 3\times 2\times 3\times 3\times 30\times 4=19440
\]
行。积分器就是十六条控制体，估计器附在同一循环里；热阻更新还读取由真值温差构成的参照，因此这些数字是一个工程核上四条策略的差别，而不是异构真值模型与在线模型配对时的感知误差——后者需要另一套积分。

档位研究另取三种食品、三种厚度、三种初温、三种含水、三种档位、四种策略和十二个种子，合计 \(3^5\times 4\times 12=11664\) 行，蒸、烤、组合按\eqref{eq:mode} 限制允许集合。两套格子的初温水平不完全相同：S1 为 \(5^\circ\mathrm{C}/20^\circ\mathrm{C}\)，档位研究为 \(4^\circ\mathrm{C}/8^\circ\mathrm{C}/20^\circ\mathrm{C}\)，因此种子 7 的演示轨迹取初温 \(8^\circ\mathrm{C}\)，并不落在 S1 的 4860 个组合里。

S0 在 30 mm、对偶、初温 \(5^\circ\mathrm{C}\)、含水 0.72、均匀场上，把每个原型的 6 个种子用 16 与 8 个控制体各积一次，18 条配对的进度差中位数约 \(0.23\)；S2 只在 30 mm、肉糜、同一初温和含水上做边界扫描，合计 480 行。二者都不并进主表分母。

\section{厚度对停机结果的主导作用}

表\ref{tab:s1}给出 S1 四个策略的总体。开环中位 \(e_{\alpha}=-0.843\)，提前比例 \(0.667\)，薄件被推到进度上沿 \(+0.150\)，厚件停在设计线之下；腔温比例中位 \(+0.150\)、能量 3310 且不停机，其上沿是进度饱和到 1 之后相对 0.85 的差。无探针提前比例 \(0.944\)，对偶 \(0.902\)：估计温度经常先于真值进度越过阈值，提前来自跟踪器把点估计送到停机门槛之前。

\begin{table}[htbp]
\centering
\caption{S1 四个策略的总体。\(e\) 为停机时 \(\alpha_m-0.85\)；\(E\) 为加热指令对时间的求和。}
\label{tab:s1}
\begin{tabular}{lrrrrr}
\toprule
策略 & \(n\) & \(e_{50}\) & \(P(e<0)\) & \(E_{50}\) & \(t_{50}/\mathrm{s}\) \\
\midrule
开环 & 4860 & \(-0.843\) & 0.667 & 1170 & 1170 \\
腔温比例 & 4860 & \(+0.150\) & 0.222 & 3310 & 3300 \\
无探针 & 4860 & \(-0.467\) & 0.944 & 1655 & 1540 \\
对偶 & 4860 & \(-0.313\) & 0.902 & 1730 & 1641 \\
\bottomrule
\end{tabular}
\end{table}

按厚度拆开之后，比较对象从策略换成几何：15 mm 开环全部停在上沿，30 mm 开环全部为负、中位 \(-0.843\)，50 mm 开环贴着 \(-0.850\)。对偶在 15 mm 中位 \(-0.417\)，30 mm \(-0.080\)，50 mm \(-0.598\) 且提前比例为 1；无探针在 30 mm 和 50 mm 上全部为负；腔温比例在 15 mm 与 30 mm 停在上沿，50 mm 中位 \(-0.461\)。三种食品的中位数几乎重合，因为停止判据读取的是冷点进度，活度只调制界面通量的幅度。

厨房里同一张面板要同时接待这三档厚度，对应三种常见进门方式。15 mm 接近一片饺子皮、一层摊开的面片，或一片鱼排：导热时间尺度在开环菜单 900 s 已经走完，Fourier 数被送到 0.80，菜单把它推过设计线，表面先湿后干，冷点来得及跟上，铃声响时进度往往已经贴着上沿。30 mm 接近一块肉饼或一枚馒头的半厚，对偶窗后方差从 0.64 落到 0.224，停机门槛跟着下降，中位最接近零，种子 7 走到 \(e_{\alpha}=+0.145\)。50 mm 接近一块带骨的厚肉或半只鸡的胸厚，同一 55 min 上限内 Fourier 数只有 0.12，探针把方差乘到 0.224，冷点仍停在设计线之下；档位研究里对偶组合的中位 \(e_{\alpha}=-0.782\)，停机 3298 s。操作者若按“中等、25 分钟”或面板上的“组合、中等、25 分钟”一刀切，给出的仍是开环 1170 s，30 mm 冷点还停在 \(16^\circ\mathrm{C}\) 量级。面板上三个档位名称改的是蒸汽和加热的许可，改不了半厚平方；厚度没有被这一次过程看见。交叉表把这一事实按厚度分列中位，而不是一句“智能菜单更准”。

说明书里的成功照片几乎总是 30 mm 这一档：表面已经上色，切开后中心不再粉色。格子把同一张照片拆开之后，15 mm 开环全部停在上沿，Fourier 数 0.80；50 mm 对偶组合中位 \(e_{\alpha}=-0.782\)，停机 3298 s，进度 0.068。照片对应的仍是种子 7、初温 \(8^\circ\mathrm{C}\)、对偶、组合那一条到达，而全表中位 \(-0.313\) 来自格子里混入的 50 mm 和 \(5^\circ\mathrm{C}\)。厨房里同一台机器同一天要接待这三档，铃声按哪一条响，决定切开之后看见的是哪一种中心。

档位研究把这一几何事实画成图\ref{fig:thick}：对偶策略下，15 mm 的中位数为正，50 mm 三种档位全部低于 \(-0.2\)；组合档 50 mm 的中位数是 \(-0.782\)，十分位到九十分位是 \([-0.819,-0.357]\)。组合档离设计线最近，50 mm 仍然全部低于设计线，说明厚度在这一批里比档位名称更决定冷点在 55 min 内能否起来。

蒸、烤、组合在面板上是三个大字，在积分器里是三条不等式。蒸档把加热夹在 0.25，对偶中位被压到 \(-0.830\)：冷凝潜热仍在，功率不够，厚件和中等件都走不完。烤档关掉蒸汽并抬高排汽，表面被推干，15 mm 和 30 mm 已经越过设计线，总体中位数因而为正，50 mm 仍然低于 \(-0.2\)。组合档放开四个量，离设计线最近，50 mm 的中位仍是 \(-0.782\)。读档位表时若先看总体中位，会把薄件的上沿误当成“烤档更熟”；按厚度拆开之后，三个字的差别落在允许集合上，冷点能不能起来仍由扩散时间决定。

\begin{figure}[htbp]
\centering
\includegraphics[width=0.88\textwidth]{F2_thickness.png}
\caption{档位研究、对偶策略：进度偏差中位数随厚度变化。\(50\,\mathrm{mm}\) 三种档位都为负。}
\label{fig:thick}
\end{figure}

图\ref{fig:d7th}把同一事实画回 S1。开环在 15 mm 为正、30 mm 与 50 mm 贴着 \(-0.84\)，因为菜单按厚度给了 900/1170/1530 s，薄件在菜单时刻已经饱和，厚件的导热时间尺度还没走完；对偶在 30 mm 最接近零，50 mm 回到 \(-0.60\)。腔温比例在 15 mm 与 30 mm 停在上沿，50 mm 才落到负半轴——它闭合的是空气温度，薄件被空气温度带着过熟，厚件的冷点跟不上腔温。图\ref{fig:d25ea}按食品名与策略给出同一张中位表，行几乎重合：停止判据读取的是冷点进度，活度只调制界面通量的幅度，三种名称换来换去，列仍由策略决定。

\begin{figure}[htbp]
\centering
\includegraphics[width=0.88\textwidth]{D7_thickness.png}
\caption{S1：停机偏差中位随厚度。对偶在 30 mm 最接近零，50 mm 仍为负。}
\label{fig:d7th}
\end{figure}

\begin{figure}[htbp]
\centering
\includegraphics[width=0.90\textwidth]{D25_e_alpha.png}
\caption{S1：食品名与策略的停机偏差中位。行几乎重合；第三列为无探针。}
\label{fig:d25ea}
\end{figure}

开环菜单终点的 Fourier 数在 15 mm 为 0.80，30 mm 为 0.26，50 mm 为 0.12，厚件在菜单时刻导热时间尺度还没有走完，与表\ref{tab:s1} 的符号一致。进度还依赖界面通量和\eqref{eq:F} 里的除数 2，Fourier 数给出的只是导热时间尺度的方向。

\section{档位变化的作用范围}

图\ref{fig:cross}是档位研究上策略与档位的交叉。开环在三档上都偏负；蒸档把加热夹在 0.25，对偶中位 \(e_{\alpha}=-0.830\)，烤档 \(+0.150\)，组合档 \(+0.146\)。烤档和组合档的总体中位数为正，是因为 15 mm 和 30 mm 已经越过设计线，50 mm 全部落在下面。蒸档的失败已经落在允许集合里：加热上限 0.25 之后，探针窗的表面温差不够，热阻更新经常不接受，进度也走不起来。档位名称改的是四个量离开零的许可，冷点在 55 min 内能否起来仍由厚度的 Fourier 数决定。对偶策略下三种食品与三档交叉时，蒸档三种食品中位都是 \(-0.83\)，烤档与组合档都是 \(+0.15\)，食品名不改变档位的符号。

\begin{figure}[htbp]
\centering
\includegraphics[width=0.88\textwidth]{F7_policy_cook.png}
\caption{档位研究：策略与档位的进度偏差中位数。开环在三档上都偏负；蒸档加热上限把其余三策压在设计线之下。}
\label{fig:cross}
\end{figure}

组合档累计加热中位数为开环 1170、对偶 1865、无探针 2319、腔温比例 3310，见图\ref{fig:energy}。烤档关闭蒸汽之后只能靠加热器，能量中位 2353，厚件的进度却不比组合档更接近设计线；蒸汽在组合档里承担了一部分近表面加热，所以同样的冷点进度用了更少的加热指令。腔温比例只看腔温、不看冷点，加热器开到时域终点。能量轴因此不是“越省越好”：开环最省，是因为它按菜单提前切断，厚件冷点还停在十几摄氏度；腔温比例最费，是因为它把空气温度固定在设定值，进度轴已经饱和仍不放手；对偶落在两者之间，是因为探针窗之后跟踪器有机会宣布到达，薄件和 30 mm 不必把加热器开满 55 min。

说明书常把能量轴单独做成卖点，说成更节能的智能菜单。开环最省，30 mm 上 1170 s 铃声响，冷点还停在 \(16^\circ\mathrm{C}\) 量级，进度偏差 \(-0.849\)，省下来的是尚未送给冷点的那一段加热。腔温比例能量中位 3310，对应 331 个记录点各乘 10 s，请求被裁在 1 上，与食材和厚度无关：空气温度被固定，进度饱和到 1 仍不放手。对偶组合档能量中位 1865，种子 7 这一条是 1864，停在 \(e_{\alpha}=+0.145\)；少花的那一截，来自停机判据开始读取这一次过程里的冷点估计，而不是来自另一套更省的传热定律。帕累托图上开环在左下，腔温比例在右上，对偶在中间偏左。把能量从 \(e_{\alpha}\) 上撕下来单独宣传，会把开环的提前切断读成优点。

\begin{figure}[htbp]
\centering
\includegraphics[width=0.88\textwidth]{F3_energy.png}
\caption{档位研究、组合档：四种策略的累计加热。箱线不含离群点。}
\label{fig:energy}
\end{figure}

S1 全表的累计加热箱线见图\ref{fig:penergy}。腔温比例压缩成一条线 3310，开环被菜单固定在 900、1170、1530 三档，无探针和对偶的箱体跨过整个时域。图\ref{fig:pstop}给出停机时刻：开环三堆对应三种厚度菜单，腔温比例堆在 55 min，无探针和对偶在中段与时域终点都有质量。能量轴的单位是加热指令对时间的求和，量纲为秒；满功率开满 55 min 对应 3300，腔温比例几乎贴着这条上沿。对偶能量中位 1730，无探针 1655：对偶略晚停、累计略高，相对腔温比例仍少加热一半。蒸汽在组合档里承担近表面加热，同一冷点进度用更少的管壁功率。

\begin{figure}[htbp]
\centering
\includegraphics[width=0.88\textwidth]{P_energy.png}
\caption{S1：累计加热箱线。纵轴是指令对时间的求和。}
\label{fig:penergy}
\end{figure}

\begin{figure}[htbp]
\centering
\includegraphics[width=0.90\textwidth]{P_stop.png}
\caption{S1：停机时刻。开环三堆是菜单；腔温比例堆在 55 min。}
\label{fig:pstop}
\end{figure}

\FloatBarrier

\section{探针的信息收益}

对偶策略 2916 条里，结束方差仍为先验 0.64 的比例是 \(1/3\)，其余落到 0.224，见图\ref{fig:var}；这个比例在 15、30、50 mm 上相同。方差是否下降，取决于新热阻是否更接近由真值温差构成的参照，与厚度无关；下降时用的乘子是 0.35，0.224 是先验乘子后的标量。厚件上的负偏差在方差下降之后仍然在：探针回答的是热阻看见了没有，这块食品在 55 min 内能否到达设计线，还要看扩散时间。

产品语言容易把这两件事合成一句“智能探针已经完成，因此可以停机”。方差降到 0.224，只说明这一次过程看见了热阻，停机门槛可以跟着下降；50 mm 的导热时间尺度仍未走完，档位研究里对偶组合的中位 \(e_{\alpha}=-0.782\)，停机 3298 s。看见热阻和走到设计线是两个判断，面板若把探针成功直接画成绿灯停机，会把厚件提前切掉。厨房里这两件事对应两盏可以分开的灯：进门后第三分钟那一阵白雾散去，方差若从 0.64 落到 0.224，第一盏灯可以亮；冷点进度走到 0.85，第二盏灯才亮。30 mm 对偶组合两盏都可以亮，种子 7 停在 1866 s，\(e_{\alpha}=+0.145\)；50 mm 对偶组合只有第一盏亮，方差已经是 0.224，进度仍停在 0.068，停机被时域终点固定在 3298 s。蒸档两盏都可以不亮：加热上限 0.25 把窗内表面温差夹住，热阻更新不接受，进度也走不起来。开环菜单按厚度档响铃，两盏都不读这一次过程里的后验。交叉表把两盏拆开写；说明书若合并成同一颗“完成”指示，厨房里看见的就是一次成功的探针和一次失败的厚件停机叠在一起。

\begin{figure}[htbp]
\centering
\includegraphics[width=0.88\textwidth]{F4_variance.png}
\caption{档位研究、对偶策略结束时的热阻方差。两堆分别是先验 0.64 和乘 \(0.35\) 之后的 0.224。}
\label{fig:var}
\end{figure}

图\ref{fig:f10ve}把结束方差对停机偏差。点落在 \(0.224\) 与 \(0.64\) 两堆。蒸档多数停在 \(0.64\)；烤档与组合档多数落到 \(0.224\)。落到 \(0.224\) 的点里仍有 \(e_{\alpha}<0\)：方差下降与进度越过设计线是两件分开的事。

\begin{figure}[htbp]
\centering
\includegraphics[width=0.88\textwidth]{F10_var_error.png}
\caption{档位研究：结束方差对停机偏差。两堆分别是先验 \(0.64\) 与乘 \(0.35\) 之后的 \(0.224\)。}
\label{fig:f10ve}
\end{figure}

种子 7、肉糜 30 mm、对偶、组合给出一条可以把中位读回过程的轨迹。方差在 192 s 由 0.64 降到 0.224，说明探针窗结束之后新热阻被接受；1680 s 仍是品质加热，冷点 \(70.21^\circ\mathrm{C}\)，\(\alpha_m=0.558\)，请求 0.476；停机在 1866 s，\(e_{\alpha}=+0.145\)，能量 1864，保持从 1872 s 开始。把 1680 s 的 \(73.0^\circ\mathrm{C}\) 与 0.755 记入这条轨迹，是八格结果，见第\ref{ch:grid}章；完整抽样见附录\ref{app:trace}。这条轨迹走到上沿，S1 对偶中位却是 \(-0.313\)，因为演示取初温 \(8^\circ\mathrm{C}\)、厚度 30 mm，S1 混有 \(5^\circ\mathrm{C}\) 和 50 mm：单条走到上沿，全表中位仍由厚度混合决定。

\section{潜热项随厚度的符号变化}

另有 216 条成对积分：三种食品、三种厚度、十二个种子，组合档、对偶策略，表面蒸汽压差打开与关闭各一次。打开之后的进度偏差减去关闭时的值，15 mm 的中位数是 \(+0.647\)，30 mm 是 \(+0.389\)，50 mm 是 \(-0.119\)；薄件上 36 对全部为正，厚件上 36 对全部为负，见图\ref{fig:antoine}。

\begin{figure}[htbp]
\centering
\includegraphics[width=0.88\textwidth]{F11_antoine.png}
\caption{打开表面蒸汽压差之后进度偏差的变化。薄件上移，厚件下移。}
\label{fig:antoine}
\end{figure}

腔温还低时，食品侧饱和压高于腔侧，差值为负，表面被吸热，厚件来不及用后面的冷凝把这口热补回来。蒸汽压公式的符号随厚度翻转；Antoine 系数取自标准水表，没有用这台机器的蒸汽表重新标定。

\section{停机偏差的总体形态}

总体中位把分布压成一个数。图\ref{fig:s1heat}按策略与厚度给出累计加热：腔温比例三档厚度都是 3310，开环随菜单上升，对偶在 15 mm 上低于无探针，在 50 mm 上两者都开到时域终点附近。图\ref{fig:boxL}把档位研究里对偶的停机偏差按厚度画成箱线，50 mm 的箱体全部在零下。

\begin{figure}[htbp]
\centering
\includegraphics[width=0.86\textwidth]{X_s1_heat_policy_thick.png}
\caption{S1：累计加热随策略与厚度。腔温比例三档重合在 3310。}
\label{fig:s1heat}
\end{figure}

\begin{figure}[htbp]
\centering
\includegraphics[width=0.88\textwidth]{P_box_thickness.png}
\caption{档位研究、对偶：停机偏差按厚度的箱线。50 mm 全部低于设计线。}
\label{fig:boxL}
\end{figure}

图\ref{fig:f5stop}给出档位研究的停止时刻，三堆对应开环菜单、对偶/无探针的中段和腔温比例的时域终点。图\ref{fig:f8t0}给出初温的边际：从冷藏升到室温，30 mm 的中位上移，内部少走大约 \(15\,\mathrm{K}\) 的焓。含水三档 \(0.64/0.72/0.80\) 的中位重合，三种食品的停止时刻也几乎重合。界面潜热在活度被裁到上沿之后对含水不再敏感，结壳又要求表面含水低于 0.45，S1 的初始含水达不到这条门槛。含水作为因子进入格子，格子却把它读成了常数。

\begin{figure}[htbp]
\centering
\includegraphics[width=0.88\textwidth]{F5_stop.png}
\caption{档位研究：停止时刻。开环落在菜单时间，腔温比例落在 3300 s。}
\label{fig:f5stop}
\end{figure}

\begin{figure}[htbp]
\centering
\includegraphics[width=0.88\textwidth]{F8_t0.png}
\caption{档位研究、对偶：初温对停机偏差。室温初值把 30 mm 中位上移。}
\label{fig:f8t0}
\end{figure}

图\ref{fig:pseedc}给出同一因子格子上 30 个种子的停机偏差。一部分饱和到上沿 \(+0.15\)，一部分停在设计线之下约 \(-0.52\)。种子号本身不翻转符号，翻转符号的是该种子落到的厚度与噪声实现。同一格子里两堆并存，说明 30 个种子已经够把“上沿或设计线之下”这一分叉固定住，不够把分叉抹平。要把分叉读成厚度而不是随机数，需要第\ref{ch:wide}章的宽格子：那里 50 mm 的上沿仍低于 \(-0.3\)，15 mm 贴着 \(+0.150\)。

\begin{figure}[htbp]
\centering
\includegraphics[width=0.88\textwidth]{P_seed_cell.png}
\caption{15 mm、初温 \(5^\circ\mathrm{C}\)、导热缺损、对偶、肉糜、含水 0.72 的 30 个种子。一簇贴上沿，一簇停在设计线之下约 \(0.525\)。}
\label{fig:pseedc}
\end{figure}

648 个格子各有 30 个种子。开环的 162 格和腔温比例的 162 格，格内 \(e_{\alpha}\) 的极差都是 0：三十个种子给出同一个停机误差，因为菜单和腔温律不读内部噪声相位。无探针 162 格的极差中位数是 0.0018，九十分位是 0.0385，最大是 0.0883，极差为 0 的格子占 0.222。对偶 162 格的极差中位数是 0.0021，九十分位是 0.0442，极差为 0 的格子占 0.333。极差大于 0.05 的格子只有 18 个，全部在 15 mm、初温 \(5^\circ\mathrm{C}\)。对偶的 9 个都是导热缺损；无探针的 9 个都是初始温度倾斜，极差 0.088，中位数 \(-0.455\)，区间 \([-0.456,-0.367]\)。三种含水的计数完全相同。

对偶、导热缺损这一档里，食材只改变两簇的人数。肉糜是 22 个种子停在 \(+0.150\)、8 个停在约 \(-0.525\)；土豆是 23 与 7；面团是 24 与 6。图\ref{fig:pseedc}只画肉糜、含水 0.72 的 30 个种子。负的一簇停在 650 s 或 652 s，能量从 975.5 到 978.1；正的一簇停在 1010 s 到 1066 s，能量从 1060.3 到 1072.3。中位数 \(+0.150\) 落在人数多的那一簇上，8 条停在设计线之下约 0.525 的轨迹不改变这个中位数。终点方差不在这些种子之间变化：对偶全部是 0.224，另外三个策略全部是 0.64。褐变终点全部是 0。

图\ref{fig:d10pol}把四个策略的 \(e_{\alpha}\) 摊开，每个策略 4860 条。零线是设计线。腔温比例的质量堆在 \(+0.15\)。图\ref{fig:pvarr}给出终点方差低于 0.3 的比例：只有对偶执行乘 0.35 的更新，另外三列是 0。对偶在三档厚度上这一比例都是 1。方差下降回答的是热阻看见了没有，进度越过设计线仍要看扩散时间。

\begin{figure}[htbp]
\centering
\includegraphics[width=0.88\textwidth]{D10_policy.png}
\caption{S1：四个策略的停机偏差分布，每个策略 4860 条。零线是设计线。}
\label{fig:d10pol}
\end{figure}

\begin{figure}[htbp]
\centering
\includegraphics[width=0.88\textwidth]{P_varr.png}
\caption{S1：终点方差低于 0.3 的比例。只有对偶走过乘 0.35 的更新。}
\label{fig:pvarr}
\end{figure}

读表时只增加图里没有写完的比较。初始含水三档在表的精度上逐格相同：对偶都是 \(e_{10}=-0.767\)、\(e_{50}=-0.313\)、提前比例 0.902、能量 1730、停机 1641。含水没有进入这一批停机误差，与活度在这三档上饱和一致。初温不改变四个策略中位数的次序，但改变幅度。对偶在 \(5^\circ\mathrm{C}\) 的中位数是 \(-0.417\)、能量 1822、停机 1740；在 \(20^\circ\mathrm{C}\) 是 \(-0.182\)、能量 1658、停机 1544。腔温比例的提前比例从 0.333 降到 0.111。厚度仍然把排序拉开得更明显。

非均质档 2 把一个控制体的导热降低，厚件的提前没有因此消失。30 mm、对偶、均匀时 \(e_{50}=-0.080\)、提前比例 1；同一厚度、对偶、初始温度倾斜 \(\pm 4^\circ\mathrm{C}\) 时 \(e_{50}=+0.036\)、提前比例降到 0.5；同一厚度、对偶、导热缺损时 \(e_{50}=-0.268\)、提前比例仍是 1。倾斜改变的是初值，不是边界驱动的冷点迁移。表中的能量单位是加热指令乘以记录间隔 10 s，不是电度。停止时刻的中位数与菜单时间在开环上重合：15 mm 为 900 s，30 mm 为 1170 s，50 mm 为 1530 s。腔温比例的能量中位数 3310 是 331 个记录点各乘 10 s，请求被裁在 1 上，与食材和厚度无关。

\subsection{交叉表的读法}
表中每一行来自 \texttt{s1\_summary.json} 的分组中位数与分位。$e$ 是停机时 $\alpha_m-0.85$。$e{<}0$ 表示停在设计线之下。$e_{50}=0.150$ 且上尾贴住 $0.150$，是进度饱和到 $1$ 之后的上沿，不是测量精度。能量 $E$ 是加热指令对时间的求和，单位为指令$\times$秒，不是电度。全部记录 \texttt{augmented=false}，\texttt{source=simulation}。这些行是同一数值核的样本，不是独立烹饪实验。
\begin{longtable}{lrrrrrrr}
\caption{四个策略的总体}\\
\toprule
分组 & $n$ & $e_{10}$ & $e_{50}$ & $e_{90}$ & $P(e{<}0)$ & $E_{50}$ & $t_{50}$ \\
\midrule
\endfirsthead
\toprule
分组 & $n$ & $e_{10}$ & $e_{50}$ & $e_{90}$ & $P(e{<}0)$ & $E_{50}$ & $t_{50}$ \\
\midrule
\endhead
对偶 & 4860 & -0.767 & -0.313 & -0.000 & 0.902 & 1730 & 1641 \\
无探针 & 4860 & -0.770 & -0.467 & -0.113 & 0.944 & 1655 & 1540 \\
开环 & 4860 & -0.850 & -0.843 & 0.150 & 0.667 & 1170 & 1170 \\
腔温比例 & 4860 & -0.698 & 0.150 & 0.150 & 0.222 & 3310 & 3300 \\
\bottomrule
\end{longtable}

\begin{longtable}{lrrrrrrr}
\caption{食材与策略}\\
\toprule
分组 & $n$ & $e_{10}$ & $e_{50}$ & $e_{90}$ & $P(e{<}0)$ & $E_{50}$ & $t_{50}$ \\
\midrule
\endfirsthead
\toprule
分组 & $n$ & $e_{10}$ & $e_{50}$ & $e_{90}$ & $P(e{<}0)$ & $E_{50}$ & $t_{50}$ \\
\midrule
\endhead
面团, 对偶 & 1620 & -0.767 & -0.313 & 0.007 & 0.900 & 1730 & 1641 \\
面团, 无探针 & 1620 & -0.770 & -0.467 & -0.113 & 0.944 & 1655 & 1540 \\
面团, 开环 & 1620 & -0.850 & -0.836 & 0.150 & 0.667 & 1170 & 1170 \\
面团, 腔温比例 & 1620 & -0.698 & 0.150 & 0.150 & 0.222 & 3310 & 3300 \\
肉糜, 对偶 & 1620 & -0.767 & -0.338 & -0.000 & 0.904 & 1730 & 1641 \\
肉糜, 无探针 & 1620 & -0.770 & -0.467 & -0.113 & 0.944 & 1655 & 1540 \\
肉糜, 开环 & 1620 & -0.850 & -0.843 & 0.150 & 0.667 & 1170 & 1170 \\
肉糜, 腔温比例 & 1620 & -0.698 & 0.150 & 0.150 & 0.222 & 3310 & 3300 \\
土豆, 对偶 & 1620 & -0.767 & -0.325 & -0.000 & 0.902 & 1729 & 1641 \\
土豆, 无探针 & 1620 & -0.770 & -0.467 & -0.113 & 0.944 & 1655 & 1540 \\
土豆, 开环 & 1620 & -0.850 & -0.836 & 0.150 & 0.667 & 1170 & 1170 \\
土豆, 腔温比例 & 1620 & -0.698 & 0.150 & 0.150 & 0.222 & 3310 & 3300 \\
\bottomrule
\end{longtable}

\begin{longtable}{lrrrrrrr}
\caption{厚度与策略}\\
\toprule
分组 & $n$ & $e_{10}$ & $e_{50}$ & $e_{90}$ & $P(e{<}0)$ & $E_{50}$ & $t_{50}$ \\
\midrule
\endfirsthead
\toprule
分组 & $n$ & $e_{10}$ & $e_{50}$ & $e_{90}$ & $P(e{<}0)$ & $E_{50}$ & $t_{50}$ \\
\midrule
\endhead
15.0, 对偶 & 1620 & -0.464 & -0.417 & 0.150 & 0.872 & 932 & 591 \\
15.0, 无探针 & 1620 & -0.478 & -0.367 & 0.150 & 0.833 & 933 & 594 \\
15.0, 开环 & 1620 & 0.150 & 0.150 & 0.150 & 0.000 & 900 & 900 \\
15.0, 腔温比例 & 1620 & 0.150 & 0.150 & 0.150 & 0.000 & 3310 & 3300 \\
30.0, 对偶 & 1620 & -0.353 & -0.080 & 0.073 & 0.833 & 1730 & 1641 \\
30.0, 无探针 & 1620 & -0.671 & -0.527 & -0.290 & 1.000 & 1655 & 1540 \\
30.0, 开环 & 1620 & -0.849 & -0.843 & -0.780 & 1.000 & 1170 & 1170 \\
30.0, 腔温比例 & 1620 & 0.150 & 0.150 & 0.150 & 0.000 & 3310 & 3300 \\
50.0, 对偶 & 1620 & -0.823 & -0.598 & -0.011 & 1.000 & 3250 & 3300 \\
50.0, 无探针 & 1620 & -0.824 & -0.627 & -0.113 & 1.000 & 3055 & 3292 \\
50.0, 开环 & 1620 & -0.850 & -0.850 & -0.850 & 1.000 & 1530 & 1530 \\
50.0, 腔温比例 & 1620 & -0.803 & -0.461 & 0.117 & 0.667 & 3310 & 3300 \\
\bottomrule
\end{longtable}

\begin{longtable}{lrrrrrrr}
\caption{食材、厚度与策略}\\
\toprule
分组 & $n$ & $e_{10}$ & $e_{50}$ & $e_{90}$ & $P(e{<}0)$ & $E_{50}$ & $t_{50}$ \\
\midrule
\endfirsthead
\toprule
分组 & $n$ & $e_{10}$ & $e_{50}$ & $e_{90}$ & $P(e{<}0)$ & $E_{50}$ & $t_{50}$ \\
\midrule
\endhead
面团, 15.0, 对偶 & 540 & -0.464 & -0.416 & 0.150 & 0.867 & 932 & 592 \\
面团, 15.0, 无探针 & 540 & -0.478 & -0.367 & 0.150 & 0.833 & 933 & 594 \\
面团, 15.0, 开环 & 540 & 0.150 & 0.150 & 0.150 & 0.000 & 900 & 900 \\
面团, 15.0, 腔温比例 & 540 & 0.150 & 0.150 & 0.150 & 0.000 & 3310 & 3300 \\
面团, 30.0, 对偶 & 540 & -0.353 & -0.080 & 0.073 & 0.833 & 1730 & 1641 \\
面团, 30.0, 无探针 & 540 & -0.689 & -0.527 & -0.290 & 1.000 & 1655 & 1540 \\
面团, 30.0, 开环 & 540 & -0.849 & -0.836 & -0.780 & 1.000 & 1170 & 1170 \\
面团, 30.0, 腔温比例 & 540 & 0.150 & 0.150 & 0.150 & 0.000 & 3310 & 3300 \\
面团, 50.0, 对偶 & 540 & -0.823 & -0.598 & -0.011 & 1.000 & 3250 & 3300 \\
面团, 50.0, 无探针 & 540 & -0.824 & -0.627 & -0.113 & 1.000 & 3053 & 3291 \\
面团, 50.0, 开环 & 540 & -0.850 & -0.850 & -0.850 & 1.000 & 1530 & 1530 \\
面团, 50.0, 腔温比例 & 540 & -0.803 & -0.461 & 0.117 & 0.667 & 3310 & 3300 \\
肉糜, 15.0, 对偶 & 540 & -0.464 & -0.417 & 0.150 & 0.878 & 932 & 591 \\
肉糜, 15.0, 无探针 & 540 & -0.478 & -0.367 & 0.150 & 0.833 & 933 & 594 \\
肉糜, 15.0, 开环 & 540 & 0.150 & 0.150 & 0.150 & 0.000 & 900 & 900 \\
肉糜, 15.0, 腔温比例 & 540 & 0.150 & 0.150 & 0.150 & 0.000 & 3310 & 3300 \\
肉糜, 30.0, 对偶 & 540 & -0.353 & -0.080 & 0.073 & 0.833 & 1730 & 1641 \\
肉糜, 30.0, 无探针 & 540 & -0.671 & -0.527 & -0.290 & 1.000 & 1655 & 1540 \\
肉糜, 30.0, 开环 & 540 & -0.850 & -0.843 & -0.821 & 1.000 & 1170 & 1170 \\
肉糜, 30.0, 腔温比例 & 540 & 0.150 & 0.150 & 0.150 & 0.000 & 3310 & 3300 \\
肉糜, 50.0, 对偶 & 540 & -0.823 & -0.598 & -0.011 & 1.000 & 3250 & 3300 \\
肉糜, 50.0, 无探针 & 540 & -0.824 & -0.627 & -0.113 & 1.000 & 3055 & 3292 \\
肉糜, 50.0, 开环 & 540 & -0.850 & -0.850 & -0.850 & 1.000 & 1530 & 1530 \\
肉糜, 50.0, 腔温比例 & 540 & -0.803 & -0.461 & 0.117 & 0.667 & 3310 & 3300 \\
土豆, 15.0, 对偶 & 540 & -0.464 & -0.417 & 0.150 & 0.872 & 932 & 591 \\
土豆, 15.0, 无探针 & 540 & -0.478 & -0.367 & 0.150 & 0.833 & 933 & 594 \\
土豆, 15.0, 开环 & 540 & 0.150 & 0.150 & 0.150 & 0.000 & 900 & 900 \\
土豆, 15.0, 腔温比例 & 540 & 0.150 & 0.150 & 0.150 & 0.000 & 3310 & 3300 \\
土豆, 30.0, 对偶 & 540 & -0.353 & -0.080 & 0.073 & 0.833 & 1729 & 1641 \\
土豆, 30.0, 无探针 & 540 & -0.671 & -0.527 & -0.290 & 1.000 & 1655 & 1540 \\
土豆, 30.0, 开环 & 540 & -0.849 & -0.836 & -0.780 & 1.000 & 1170 & 1170 \\
土豆, 30.0, 腔温比例 & 540 & 0.150 & 0.150 & 0.150 & 0.000 & 3310 & 3300 \\
土豆, 50.0, 对偶 & 540 & -0.823 & -0.598 & -0.011 & 1.000 & 3249 & 3300 \\
土豆, 50.0, 无探针 & 540 & -0.824 & -0.627 & -0.113 & 1.000 & 3055 & 3292 \\
土豆, 50.0, 开环 & 540 & -0.850 & -0.850 & -0.850 & 1.000 & 1530 & 1530 \\
土豆, 50.0, 腔温比例 & 540 & -0.803 & -0.461 & 0.117 & 0.667 & 3310 & 3300 \\
\bottomrule
\end{longtable}

\begin{longtable}{lrrrrrrr}
\caption{非均质档、厚度与策略}\\
\toprule
分组 & $n$ & $e_{10}$ & $e_{50}$ & $e_{90}$ & $P(e{<}0)$ & $E_{50}$ & $t_{50}$ \\
\midrule
\endfirsthead
\toprule
分组 & $n$ & $e_{10}$ & $e_{50}$ & $e_{90}$ & $P(e{<}0)$ & $E_{50}$ & $t_{50}$ \\
\midrule
\endhead
0, 15.0, 对偶 & 540 & -0.444 & -0.391 & -0.338 & 1.000 & 924 & 578 \\
0, 15.0, 无探针 & 540 & -0.479 & -0.422 & -0.365 & 1.000 & 925 & 580 \\
0, 15.0, 开环 & 540 & 0.150 & 0.150 & 0.150 & 0.000 & 900 & 900 \\
0, 15.0, 腔温比例 & 540 & 0.150 & 0.150 & 0.150 & 0.000 & 3310 & 3300 \\
0, 30.0, 对偶 & 540 & -0.155 & -0.080 & -0.013 & 1.000 & 1714 & 1617 \\
0, 30.0, 无探针 & 540 & -0.579 & -0.513 & -0.410 & 1.000 & 1640 & 1516 \\
0, 30.0, 开环 & 540 & -0.849 & -0.835 & -0.788 & 1.000 & 1170 & 1170 \\
0, 30.0, 腔温比例 & 540 & 0.150 & 0.150 & 0.150 & 0.000 & 3310 & 3300 \\
0, 50.0, 对偶 & 540 & -0.767 & -0.409 & -0.052 & 1.000 & 3229 & 3300 \\
0, 50.0, 无探针 & 540 & -0.770 & -0.529 & -0.268 & 1.000 & 3024 & 3249 \\
0, 50.0, 开环 & 540 & -0.850 & -0.850 & -0.850 & 1.000 & 1530 & 1530 \\
0, 50.0, 腔温比例 & 540 & -0.698 & -0.299 & 0.100 & 0.500 & 3310 & 3300 \\
1, 15.0, 对偶 & 540 & -0.418 & -0.362 & -0.305 & 1.000 & 931 & 583 \\
1, 15.0, 无探针 & 540 & -0.456 & -0.352 & -0.336 & 1.000 & 932 & 584 \\
1, 15.0, 开环 & 540 & 0.150 & 0.150 & 0.150 & 0.000 & 900 & 900 \\
1, 15.0, 腔温比例 & 540 & 0.150 & 0.150 & 0.150 & 0.000 & 3310 & 3300 \\
1, 30.0, 对偶 & 540 & -0.001 & 0.036 & 0.095 & 0.500 & 1741 & 1643 \\
1, 30.0, 无探针 & 540 & -0.489 & -0.390 & -0.289 & 1.000 & 1664 & 1539 \\
1, 30.0, 开环 & 540 & -0.849 & -0.834 & -0.780 & 1.000 & 1170 & 1170 \\
1, 30.0, 腔温比例 & 540 & 0.150 & 0.150 & 0.150 & 0.000 & 3310 & 3300 \\
1, 50.0, 对偶 & 540 & -0.756 & -0.383 & -0.011 & 1.000 & 3255 & 3300 \\
1, 50.0, 无探针 & 540 & -0.759 & -0.436 & -0.093 & 1.000 & 3071 & 3281 \\
1, 50.0, 开环 & 540 & -0.850 & -0.850 & -0.850 & 1.000 & 1530 & 1530 \\
1, 50.0, 腔温比例 & 540 & -0.678 & -0.280 & 0.117 & 0.500 & 3310 & 3300 \\
2, 15.0, 对偶 & 540 & -0.525 & -0.463 & 0.150 & 0.617 & 941 & 607 \\
2, 15.0, 无探针 & 540 & -0.412 & -0.130 & 0.150 & 0.500 & 988 & 799 \\
2, 15.0, 开环 & 540 & 0.150 & 0.150 & 0.150 & 0.000 & 900 & 900 \\
2, 15.0, 腔温比例 & 540 & 0.150 & 0.150 & 0.150 & 0.000 & 3310 & 3300 \\
2, 30.0, 对偶 & 540 & -0.353 & -0.268 & -0.182 & 1.000 & 1749 & 1670 \\
2, 30.0, 无探针 & 540 & -0.689 & -0.618 & -0.564 & 1.000 & 1674 & 1568 \\
2, 30.0, 开环 & 540 & -0.850 & -0.844 & -0.824 & 1.000 & 1170 & 1170 \\
2, 30.0, 腔温比例 & 540 & 0.150 & 0.150 & 0.150 & 0.000 & 3310 & 3300 \\
2, 50.0, 对偶 & 540 & -0.823 & -0.631 & -0.439 & 1.000 & 3245 & 3300 \\
2, 50.0, 无探针 & 540 & -0.824 & -0.659 & -0.479 & 1.000 & 3066 & 3292 \\
2, 50.0, 开环 & 540 & -0.850 & -0.850 & -0.850 & 1.000 & 1530 & 1530 \\
2, 50.0, 腔温比例 & 540 & -0.803 & -0.523 & -0.244 & 1.000 & 3310 & 3300 \\
\bottomrule
\end{longtable}

\begin{longtable}{lrrrrrrr}
\caption{初始含水与策略}\\
\toprule
分组 & $n$ & $e_{10}$ & $e_{50}$ & $e_{90}$ & $P(e{<}0)$ & $E_{50}$ & $t_{50}$ \\
\midrule
\endfirsthead
\toprule
分组 & $n$ & $e_{10}$ & $e_{50}$ & $e_{90}$ & $P(e{<}0)$ & $E_{50}$ & $t_{50}$ \\
\midrule
\endhead
0.64, 对偶 & 1620 & -0.767 & -0.313 & -0.000 & 0.902 & 1730 & 1641 \\
0.64, 无探针 & 1620 & -0.770 & -0.467 & -0.113 & 0.944 & 1655 & 1540 \\
0.64, 开环 & 1620 & -0.850 & -0.843 & 0.150 & 0.667 & 1170 & 1170 \\
0.64, 腔温比例 & 1620 & -0.698 & 0.150 & 0.150 & 0.222 & 3310 & 3300 \\
0.72, 对偶 & 1620 & -0.767 & -0.313 & -0.000 & 0.902 & 1730 & 1641 \\
0.72, 无探针 & 1620 & -0.770 & -0.467 & -0.113 & 0.944 & 1655 & 1540 \\
0.72, 开环 & 1620 & -0.850 & -0.843 & 0.150 & 0.667 & 1170 & 1170 \\
0.72, 腔温比例 & 1620 & -0.698 & 0.150 & 0.150 & 0.222 & 3310 & 3300 \\
0.8, 对偶 & 1620 & -0.767 & -0.313 & -0.000 & 0.902 & 1730 & 1641 \\
0.8, 无探针 & 1620 & -0.770 & -0.467 & -0.113 & 0.944 & 1655 & 1540 \\
0.8, 开环 & 1620 & -0.850 & -0.843 & 0.150 & 0.667 & 1170 & 1170 \\
0.8, 腔温比例 & 1620 & -0.698 & 0.150 & 0.150 & 0.222 & 3310 & 3300 \\
\bottomrule
\end{longtable}

\begin{longtable}{lrrrrrrr}
\caption{初始温度与策略}\\
\toprule
分组 & $n$ & $e_{10}$ & $e_{50}$ & $e_{90}$ & $P(e{<}0)$ & $E_{50}$ & $t_{50}$ \\
\midrule
\endfirsthead
\toprule
分组 & $n$ & $e_{10}$ & $e_{50}$ & $e_{90}$ & $P(e{<}0)$ & $E_{50}$ & $t_{50}$ \\
\midrule
\endhead
5.0, 对偶 & 2430 & -0.823 & -0.417 & -0.000 & 0.915 & 1822 & 1740 \\
5.0, 无探针 & 2430 & -0.824 & -0.579 & 0.150 & 0.889 & 1738 & 1628 \\
5.0, 开环 & 2430 & -0.850 & -0.849 & 0.150 & 0.667 & 1170 & 1170 \\
5.0, 腔温比例 & 2430 & -0.803 & 0.150 & 0.150 & 0.333 & 3310 & 3300 \\
20.0, 对偶 & 2430 & -0.463 & -0.182 & 0.073 & 0.889 & 1658 & 1544 \\
20.0, 无探针 & 2430 & -0.564 & -0.366 & -0.113 & 1.000 & 1587 & 1448 \\
20.0, 开环 & 2430 & -0.850 & -0.821 & 0.150 & 0.667 & 1170 & 1170 \\
20.0, 腔温比例 & 2430 & -0.244 & 0.150 & 0.150 & 0.111 & 3310 & 3300 \\
\bottomrule
\end{longtable}


\section{能量、方差与褐变的伴随量}
下列数字由 \texttt{s1\_summary.json} 的 19440 行重新分组。能量仍是加热指令乘记录间隔。方差低于 0.3 表示走过一次乘 0.35 的更新；未更新的先验方差是 0.64。褐变是终点 $\alpha_b$，不是停机时刻。
\begin{longtable}{lrrrrrr}
\caption{按厚度与策略的能量分位与终点方差}\\
\toprule
分组 & $n$ & $E_{10}$ & $E_{50}$ & $E_{90}$ & 方差$<0.3$ & $\alpha_{b,50}$ \\
\midrule\endfirsthead
\toprule 分组 & $n$ & $E_{10}$ & $E_{50}$ & $E_{90}$ & 方差$<0.3$ & $\alpha_{b,50}$ \\\midrule\endhead
15\,mm, 开环 & 1620 & 900 & 900 & 900 & 0.000 & 0.0000 \\
15\,mm, 腔温比例 & 1620 & 3310 & 3310 & 3310 & 0.000 & 0.0000 \\
15\,mm, 无探针 & 1620 & 892 & 933 & 1076 & 0.000 & 0.0000 \\
15\,mm, 对偶 & 1620 & 890 & 932 & 1064 & 1.000 & 0.0000 \\
30\,mm, 开环 & 1620 & 1170 & 1170 & 1170 & 0.000 & 0.0000 \\
30\,mm, 腔温比例 & 1620 & 3310 & 3310 & 3310 & 0.000 & 0.0000 \\
30\,mm, 无探针 & 1620 & 1563 & 1655 & 1753 & 0.000 & 0.0000 \\
30\,mm, 对偶 & 1620 & 1629 & 1730 & 1836 & 1.000 & 0.0000 \\
50\,mm, 开环 & 1620 & 1530 & 1530 & 1530 & 0.000 & 0.0000 \\
50\,mm, 腔温比例 & 1620 & 3310 & 3310 & 3310 & 0.000 & 0.0000 \\
50\,mm, 无探针 & 1620 & 2901 & 3055 & 3178 & 0.000 & 0.0000 \\
50\,mm, 对偶 & 1620 & 3157 & 3250 & 3307 & 1.000 & 0.0000 \\
\bottomrule\end{longtable}
对偶在三档厚度上的方差下降比例都是 1。开环、腔温比例和无探针的比例是 0，因为它们的代码路径不写 $\hat r$。

终点褐变在这 19440 行里不是“中位数碰巧为零”。$\alpha_b$ 的最小值、十分位、中位数、九十分位和最大值都是 0，大于零的比例是 0。褐变增量要求表面含水低于 0.55，同时表面温度高于 $110^\circ\mathrm{C}$。S1 的初始含水是 0.64、0.72 和 0.80。表面失水项再乘上 $2\times 10^{-9}\,\mathrm{m}^2/\mathrm{s}$ 量级的扩散系数之后，下降率约为 $10^{-11}\,\mathrm{s}^{-1}$，3300 s 的下降量级是 $10^{-8}$。含水到不了 0.55，褐变增量在每一步都是 0，不是四位小数把一个小正数抹掉。结壳门限是 0.45，同样到不了，所以路径 C 里结壳增大热阻的那一支在这次战役里没有被积分。三个原型的活度在这三档含水上都贴着 0.99，表里的食材列因此也不能比较蛋白与淀粉的表面反应。

能量列里有一条与电功率无关的恒等式。战役记录间隔是 $2\,\mathrm{s}\times 5=10\,\mathrm{s}$，55 min 共 1650 步，记录点是 331 个。腔温比例把请求裁到 1，加热器平衡在 $196.5^\circ\mathrm{C}$，削减不动作，每个记录点的指令都是 1，能量就是 $331\times 10=3310$。开环在菜单终点把指令写成 0，900 s、1170 s、1530 s 都落在 10 s 网格上，能量中位数就等于菜单秒数。末步若不是记录步，代码仍把它记入，并仍乘满 10 s；3310 比 3300 多出的那一格是这个记账。


档位研究里 50 mm 的对偶样本仍全部低于设计线，薄件把总体中位抬了上去，所以先按厚度拆开再比较策略。结束方差 \(0.224\) 是先验 0.64 乘 0.35：探针回答的是热阻看见了没有，进度越过设计线仍要看扩散时间。19440 与 11664 是两套格子的样本条数，分母分开写。种子 7 把交叉表读回一条可见轨迹：方差在 192 s 由 0.64 降到 0.224，1866 s 停机，\(e_{\alpha}=+0.145\)，能量 1864；15 mm 同一策略把进度推到上沿，50 mm 停在 \(e_{\alpha}=-0.782\)，三条结束方差都是 0.224。网格、消融和轨迹族把这一符号固定得更死：八格会翻号，含水三档给出零效应，50 mm 要把时域放到 70 min 才走到同一条设计线。

说明书里的成功照片对应格子里人数多的那一簇。15 mm、初温 \(5^\circ\mathrm{C}\)、导热缺损、对偶这一格里，肉糜 22 个种子停在 \(+0.150\)，8 个停在约 \(-0.525\)，中位数落在人数多的簇上；厨房里拍到的那一张切开图，往往来自已经饱和的薄件或已经回温的 30 mm，50 mm 在档位研究里对偶组合全部低于设计线，拍不到这张照片。交叉表把照片拆开之后，厚度符号先于策略名称。演示页默认读的是 30 mm、\(8^\circ\mathrm{C}\) 那一条成功轨迹，停机显示 1864 s、\(e_{\alpha}=+0.145\)；切到 50 mm，同一策略停在 3298 s、\(-0.782\)，两张图都来自已经写出的数组。S1 全表的能量中位是开环 1170、对偶 1730、无探针 1655、腔温比例 3310，单位是加热指令乘以记录间隔 10 s，不是电度；把能量从 \(e_{\alpha}\) 上撕下来单独宣传，会把开环的提前切断读成优点。

\FloatBarrier
